\subsection{A distinct frozen candidate at weight nine}
\label{s8new:sec:candidate}
The later signed/zero continuation supplies a new $S_8$ candidate in a
Gaussian basket with odd inner indices. It is distinct from the vector
rigorously rejected in Proposition~\ref{research:prop:S8-rejected}.

\begin{conjecture}[The next mixed Gaussian reduction]\label{s8new:conj:S8}
\begin{align}
 S_8={}&\frac{52334}{37973}g_{8,1}+\frac{2824}{37973}g_{6,3}
       -\frac{6832}{189865}g_{4,5}-\frac{139760}{265811}g_{2,7}\notag\\
 &+\frac{998183\pi^9}{17283022848}
       -\frac{78615}{19442176}G\zeta(7)
       -\frac{201021}{4252976}\beta(4)\zeta(5)\notag\\
 &-\frac{290505}{1063244}\beta(6)\zeta(3)-2\beta(8)\log2.
 \label{s8new:eq:formula}
\end{align}
\end{conjecture}
In the order
\[
 (S_8,g_{8,1},g_{6,3},g_{4,5},g_{2,7},\pi^9,
 G\zeta(7),\beta(4)\zeta(5),\beta(6)\zeta(3),\beta(8)\log2),
\]
its primitive integer vector is
\[
\begin{split}
 (&-10974719508480,15125246115840,816174858240,-394908401664,\\
 &-5770366156800,633846205,-44376595200,-518730670080,\\
 &-2998569369600,-21949439016960).
\end{split}
\]
All entries have greatest common divisor one; each displayed right-side
coefficient is minus its vector entry divided by the first entry.

The discovery used 720 Euler terms at 250 working digits and then PSLQ
at 150 digits, with tolerance $10^{-140}$ and coefficient bound $10^{20}$.
After freezing the vector, separate Mellin integrals at 220 digits and a
1400-term Euler evaluation at 500 digits corroborated it. Their residuals
are numerical diagnostics and are not used to prove equality.

The fresh standard-library replay reconstructs two rational proximity
certificates: the 1100-term source-bound scheme and the 1200-term direct
scheme. For $R_8$ equal to the left side minus the right side of
\eqref{s8new:eq:formula}, the latter proves
\[
 |R_8|<10^{-355}.
\]
The same scheme improves the normalized $S_6$ residual to that bound.
Each residual interval contains zero. The exact endpoints and frozen
vectors are preserved with the signed/zero continuation. These enclosures
prove proximity and leave both identities conjectural.
